\section{Why Responsible AI}

\begin{frame}{Where This Lecture Sits}

    So far the course has optimised a single number: accuracy, F1, RMSE. That number is computed
    \textbf{over the whole test set}.

    \vspace{1em}

    \begin{block}{The gap this lecture closes}
        A model that is accurate \emph{on average} can still be systematically wrong for a
        particular group of people, and can be right for reasons you would not accept if you
        could see them.
    \end{block}

    \vspace{1em}

    \begin{itemize}
        \setlength\itemsep{0.5em}
        \item \textbf{Part 1 --- Fairness:} how to find and reduce systematic error across groups.
        \item \textbf{Part 2 --- Interpretability:} how to find out what a model is actually using.
        \item \textbf{Part 3 --- Documentation and governance:} how to write it down and who checks it.
    \end{itemize}

    \vspace{0.5em}
    \centering
    \small
    \textit{Not covered here:} agent- and LLM-specific safety --- prompt injection, tool misuse,
    excessive agency --- which is the subject of \textbf{Lecture 12 of the NLP course}.

    \vspace{0.6em}
    \normalsize
    \textit{We start with four documented cases. Each one failed in a different place in the pipeline.}
\end{frame}

% ---------------------------------------------------------------------------
% Case 1: COMPAS
% ---------------------------------------------------------------------------
\begin{frame}{Case 1: COMPAS Recidivism Scores (2016)}

    \textbf{The system.} COMPAS is a commercial risk-assessment tool (Northpointe, now Equivant)
    that scores a criminal defendant's likelihood of re-offending on a 1--10 scale. US courts have
    used such scores to inform bail, sentencing, and parole decisions.

    \vspace{0.8em}

    \textbf{The analysis.} ProPublica journalists obtained scores for more than 7{,}000 people
    arrested in Broward County, Florida in 2013--2014 and checked, two years later, who actually
    re-offended.

    \vspace{0.8em}

    \begin{columns}[T]
        \begin{column}{0.52\textwidth}
            \centering
            \small
            \begin{tabular}{p{3.5cm}cc}
                \toprule
                & \textbf{White} & \textbf{Black} \\
                \midrule
                Labelled higher risk, but did \emph{not} re-offend & 23.5\% & 44.9\% \\
                \addlinespace
                Labelled lower risk, yet \emph{did} re-offend & 47.7\% & 28.0\% \\
                \bottomrule
            \end{tabular}
        \end{column}
        \begin{column}{0.44\textwidth}
            \begin{itemize}\small
                \item Overall, 61\% of those scored likely to re-offend were arrested for some
                      crime within two years.
                \item Of those predicted to commit a \emph{violent} crime, 20\% did.
            \end{itemize}
        \end{column}
    \end{columns}

    \vspace{0.5em}
    {\scriptsize J.~Angwin, J.~Larson, S.~Mattu and L.~Kirchner, ``Machine Bias,'' \emph{ProPublica},
    23 May 2016, \href{https://www.propublica.org/article/machine-bias-risk-assessments-in-criminal-sentencing}{propublica.org}.}
\end{frame}

\begin{frame}{Case 1: The Disagreement That Followed}

    Northpointe replied that its score was \textbf{equally accurate for both groups}: among
    defendants given the same score, roughly the same fraction re-offended regardless of race.

    \vspace{0.8em}

    \begin{block}{Both sides were reporting true numbers}
        ProPublica measured \textbf{error rates} conditioned on the true outcome
        (who got wrongly flagged). Northpointe measured \textbf{calibration} conditioned on the
        score (what a score means). These are different quantities.
    \end{block}

    \vspace{0.8em}

    \begin{itemize}
        \setlength\itemsep{0.4em}
        \item Chouldechova (2017) and Kleinberg, Mullainathan and Raghavan (2016) showed
              independently that when the two groups have \textbf{different base rates} of the
              outcome, a score cannot satisfy both criteria at once, except in degenerate cases.
        \item So this was not a bug that better engineering would fix. It was a
              \textbf{choice between incompatible definitions of fair} --- and nobody had made it
              explicitly.
    \end{itemize}

    \vspace{0.5em}
    \centering
    \textit{We return to this result in Section 4.}
\end{frame}

% ---------------------------------------------------------------------------
% Case 2: Amazon hiring tool
% ---------------------------------------------------------------------------
\begin{frame}[allowframebreaks]{Case 2: A Resume-Screening Tool (Reuters reporting, 2018)}

    \begin{tcolorbox}[colback=raiAmber!5, colframe=raiAmber!70, title=Note on evidence]
        \small This case comes from \textbf{news reporting based on anonymous sources}, not from a
        published study. Treat the mechanism as illustrative and the details as journalistic.
    \end{tcolorbox}

    \vspace{0.6em}

    According to Reuters, a team at Amazon built experimental tools from 2014 that scored job
    applicants one to five stars. By 2015 the team found the system was not rating candidates for
    software-developer roles in a gender-neutral way.

    \vspace{0.6em}

    \begin{itemize}
        \setlength\itemsep{0.4em}
        \item The models were trained on resumes submitted to the company over roughly a decade ---
              a pool that was predominantly male in technical roles.
        \item The reporting states the system penalised resumes containing the word
              ``women's'' and downgraded graduates of two all-women's colleges.
        \item Amazon said the tool was never used by recruiters to evaluate candidates, and the
              project was abandoned.
    \end{itemize}

    \vspace{0.5em}

    \begin{block}{The mechanism, stated plainly}
        \small The label was not ``a good engineer.'' It was ``a person this company hired
        before.'' Removing the gender field does not help, because the model finds
        \textbf{proxies}. The learning worked; the \textbf{target definition} was the defect.
    \end{block}

    \vspace{0.3em}
    {\scriptsize J.~Dastin, ``Amazon scraps secret AI recruiting tool that showed bias against
    women,'' \emph{Reuters}, 10 October 2018.}
\end{frame}

% ---------------------------------------------------------------------------
% Case 3: Gender Shades
% ---------------------------------------------------------------------------
\begin{frame}{Case 3: Gender Shades --- Disaggregated Error Rates (2018)}

    Buolamwini and Gebru audited three \textbf{commercial} gender-classification APIs on the
    \emph{Pilot Parliaments Benchmark} (PPB): 1{,}270 parliamentarians from three African and three
    European countries, labelled by Fitzpatrick skin type and by binary gender.

    \vspace{0.18em}
    \centering
    \scriptsize
    \begin{tabular}{lccccc}
        \toprule
        \textbf{Classifier} & \textbf{All} & \textbf{Darker female} & \textbf{Darker male}
        & \textbf{Lighter female} & \textbf{Lighter male} \\
        \midrule
        Microsoft & 6.3\%  & 20.8\% & 6.0\%  & 1.7\% & 0.0\% \\
        Face++    & 10.0\% & 34.5\% & 0.7\%  & 9.8\% & 0.8\% \\
        IBM       & 12.1\% & 34.7\% & 12.0\% & 7.1\% & 0.3\% \\
        \bottomrule
    \end{tabular}

    \vspace{0.12em}
    {\scriptsize Error rate $=1-$ true positive rate, evaluated April--May 2017.}

    \vspace{0.18em}
    \raggedright
    \begin{columns}[T]
        \begin{column}{0.51\textwidth}
            \begin{itemize}\scriptsize
                \setlength\itemsep{0.1em}
                \item The aggregate column looks acceptable. The intersectional columns do not.
                \item Marginal slices hide it: an 8.1\%--20.6\% gap by gender alone, 11.8\%--19.2\%
                      by skin type alone.
                \item Maximum reported error-rate gap between the best- and worst-classified
                      subgroups: \textbf{34.4\%}.
            \end{itemize}
        \end{column}
        \begin{column}{0.47\textwidth}
            \begin{tcolorbox}[colback=raiBlue!5, colframe=raiBlue!60, boxrule=0.6pt]
                \scriptsize
                \textbf{Two lessons.} (1) Slice by every group \emph{and by combinations}, and
                report the worst slice with its sample size. (2) Quote the study for what it
                measured --- binary gender classification, one benchmark, three APIs, one point in
                time --- not as ``face recognition does not work.''
            \end{tcolorbox}
        \end{column}
    \end{columns}

    \vspace{0.12em}
    {\scriptsize J.~Buolamwini and T.~Gebru, ``Gender Shades: Intersectional Accuracy Disparities in
    Commercial Gender Classification,'' \emph{PMLR} 81:77--91, FAT* 2018, Table 4 and Section 4.1.}
\end{frame}

% ---------------------------------------------------------------------------
% Case 4: Obermeyer
% ---------------------------------------------------------------------------
\begin{frame}[allowframebreaks]{Case 4: A Health Risk Algorithm and Its Proxy Label (2019)}

    \textbf{The system.} A widely deployed algorithm identified patients who should be enrolled in
    high-intensity care-management programmes. It predicted \textbf{future health-care costs} as a
    stand-in for \textbf{future health needs}.

    \vspace{0.8em}

    \begin{block}{The defect}
        Less money is spent on Black patients at the same level of illness. Predicting cost
        therefore predicts \emph{spending}, not \emph{sickness} --- and the algorithm concluded that
        equally sick Black patients were healthier.
    \end{block}

    \vspace{0.8em}

    \begin{itemize}
        \setlength\itemsep{0.4em}
        \item At a given risk score, Black patients were measurably sicker (more uncontrolled
              chronic conditions).
        \item Reformulating the label removed most of the disparity: the share of Black patients
              flagged for extra help rose from \textbf{17.7\% to 46.5\%}.
        \item The sensitive attribute was \textbf{not} an input. The \textbf{label} carried the bias.
    \end{itemize}

    \vspace{0.4em}
    {\scriptsize Z.~Obermeyer, B.~Powers, C.~Vogeli and S.~Mullainathan, ``Dissecting racial bias in
    an algorithm used to manage the health of populations,'' \emph{Science} 366(6464):447--453, 2019.}
\end{frame}

\begin{frame}{What the Four Cases Have in Common}

    \centering
    \small
    \begin{tabular}{p{2.6cm} p{4.4cm} p{4.6cm}}
        \toprule
        \textbf{Case} & \textbf{Where it broke} & \textbf{What would have caught it} \\
        \midrule
        COMPAS & Conflicting fairness criteria, never chosen explicitly
               & Naming the fairness criterion before deployment \\
        \addlinespace
        Resume screening & Target definition (``past hires'') and proxies
               & Asking what the label actually measures \\
        \addlinespace
        Gender Shades & Benchmark composition and unreported slices
               & Disaggregated, intersectional evaluation \\
        \addlinespace
        Health risk score & Proxy label (cost for need)
               & Auditing the label against the real outcome \\
        \bottomrule
    \end{tabular}

    \vspace{1.2em}

    \begin{block}{}
        \centering
        None of these were caused by a bad optimiser or a small test set. All four were
        \textbf{specification} problems that a standard evaluation report would not surface.
    \end{block}
\end{frame}
